\section{Methods}
\label{sec:methods}

\subsection{Regions and locations}
\label{sec:methods-regions}

Every forecast is scored globally, for three regions of
similar size (Table~\ref{tab:regions}) and for five metropolitan areas,
each represented by the $3\times3$ grid points
($0.75^\circ\times0.75^\circ$) nearest to its centre. The regions cover the
mid-latitudes of two continents and a monsoon region; North America is the
WMO's standard verification area \citep{wmo2010gdpfs}. The
locations span oceanic, monsoon, Mediterranean, continental and
subtropical-highland climates in both hemispheres. The Bay Area box
includes points over the bay and the Pacific. In addition, error sums on
$10^\circ$ cells give latitude bands
(\S\ref{sec:methods-further}).

\begin{table}[tbp]
\centering
\small
\caption{Regions and locations used for verification. Locations: the
$3\times3$ grid points nearest to the given centre.}
\label{tab:regions}
\begin{tabular}{lll}
\toprule
Name & Area & Climate \\
\midrule
UK/Europe & 35--72$^\circ$N, 15$^\circ$W--40$^\circ$E & mid-latitude \\
North America & 25--60$^\circ$N, 145--50$^\circ$W & mid-latitude, continental \\
South Asia & 5--35$^\circ$N, 60--100$^\circ$E & tropical and subtropical monsoon \\
\midrule
London & 51.5$^\circ$N, 0.1$^\circ$W & temperate oceanic \\
Delhi & 28.6$^\circ$N, 77.2$^\circ$E & semi-arid, monsoon \\
Bay Area & 37.6$^\circ$N, 122.2$^\circ$W & Mediterranean, coastal \\
Beijing & 39.9$^\circ$N, 116.4$^\circ$E & continental monsoon \\
S\~{a}o Paulo & 23.55$^\circ$S, 46.65$^\circ$W & subtropical highland \\
\bottomrule
\end{tabular}
\end{table}

\subsection{Paired fp32 and fp16 forecasts}
\label{sec:methods-tier1}

For each model and initial date, we run the same 10-day forecast from
identical initial conditions with the network in fp32 and in fp16
(\S\ref{sec:quantisation}), and score both against ERA5
(\S\ref{sec:verification-data}) at daily lead times (days 1--10). Because
the two forecasts share the initial state and the weather situation,
their difference in RMSE isolates the effect of precision from the much
larger variation of forecast error between dates. We also record the RMS
difference between the fp16 and the fp32 forecast, their
\emph{departure}, which measures how much the precision changes the
forecast regardless of whether it becomes more or less accurate. The bf16
forecasts used for the secondary comparison (\S\ref{sec:results-bf16})
are built and scored in the same way.

The initial dates are drawn from 2024 by stratified random sampling:
one date from each of the 52 calendar weeks, with the initialisation
times 00, 06, 12 and 18\,UTC used 13 times each and the days of the week
used as evenly as possible. All models use the same 52
initialisations, so that differences between them cannot arise from
differences in the weather situations they were tested on. Initialisations about a week
apart give forecasts whose errors are close to independent between
dates. The sensitivity experiments and the ensemble
(\S\ref{sec:methods-tier2}, \S\ref{sec:methods-further}) use every fourth
of these dates (13).

For each model, variable, lead time and area, we report the change in
mean RMSE from fp32 to fp16 relative to the fp32 RMSE, with a 95\%
confidence interval and a two-sided paired $t$-test of the per-date RMSE
differences at the 5\% level. With 50 combinations of variable and lead
time per model and area, we do not correct for multiple testing, and
treat isolated significant results, especially at the locations, with
caution. To judge whether a change matters in practice, we also express
it as the lead time over which the fp32 RMSE grows by the same amount
(\S\ref{sec:results-fp16}).

To show what this design can detect, we compute for each combination
the minimum detectable effect \citep{bloom1995minimum,cohen1988statistical}:
the smallest change in mean RMSE that the paired $t$-test detects with
80\% probability, given $N=52$ dates and the observed standard deviation
$s$ of the per-date differences, $(t_{0.975,N-1}+t_{0.8,N-1})\,s/\sqrt{N}$
(Appendix~\ref{sec:appendix-mde}, Table~\ref{tab:power}). This is based on the observed variability, not
on the observed effect, so it describes the sensitivity of the design
rather than the probability that the reported effects are real.
Globally, the median detectable change is 0.03--0.33\% of the fp32 RMSE,
depending on the model; over the three regions, where scores vary more
from date to date, 0.08--1.1\%. Because the fp16 and fp32 forecasts are
so closely paired, the design resolves changes far smaller than any that
would matter in practice: a significant result here shows that fp16
changes the forecast, not that the change is important.

\begin{table}[htbp]
\centering
\small
\caption{Smallest change in mean RMSE, as \% of the fp32 RMSE, that the paired $t$-test detects with probability 0.8 at the 5\% level, given the $N$ dates and the observed date-to-date variability of the fp16$-$fp32 differences. Median over the 50 combinations of variable and lead time per area (global: with range).}
\label{tab:power}
\begin{tabular}{lrrrrr}
\toprule
Model & $N$ & Global & UK/Europe & N.\ America & S.\ Asia \\
\midrule
FastNet & 52 & 0.33 [0.01, 0.7] & 1.11 & 1.00 & 0.52 \\
Pangu-Weather & 52 & 0.05 [0.01, 0.2] & 0.16 & 0.13 & 0.14 \\
FourCastNet & 52 & 0.23 [0.02, 6.5] & 0.95 & 0.66 & 1.13 \\
SFNO & 52 & 0.03 [0.00, 0.1] & 0.11 & 0.08 & 0.09 \\
Aurora & 52 & 0.12 [0.01, 0.5] & 0.36 & 0.30 & 0.54 \\
\bottomrule
\end{tabular}
\end{table}


\subsection{Sensitivity experiments}
\label{sec:methods-tier2}

To find where in the network fp16 rounding enters the forecast, we
use a one-at-a-time (OAT) design on the components listed in
Table~\ref{tab:components}: one component at a time runs in fp16 and the
rest of the network in fp32, which measures the departure that the
component causes on its own. The inputs of the component are rounded to
fp16 when they enter it, and its outputs converted back to fp32 when they
leave it. This corresponds to the noising variant of activation patching, in which a
component is perturbed in an otherwise clean model
\citep{heimersheim2024activation}, a technique used to locate computations
in language models \citep{meng2022locating}. One-at-a-time designs cannot detect
interactions between components \citep{saltelli2010perfunctory}; we
check for them by comparing the components' effects with that of the
whole network (\S\ref{sec:results-sensitivity}), and we also run the
network with every component converted in this way.

Each forecast is compared with the fp32 forecast from the same
initial conditions rather than with ERA5: the departure, the
area-weighted RMS difference between the two, is computed globally for
each variable at every model step up to 24\,h and averaged over the 13
dates. A comparison with ERA5 would add the much larger forecast error
that both forecasts share. We express the departure $D_c$ with component
$c$ in fp16 relative to the departure $D_\mathrm{fp16}$ of the whole
network in fp16, as the share $D_c/D_\mathrm{fp16}$ of the departure that
the component causes on its own.

The one-at-a-time design asks where in the network the rounding
matters. Practice in other applications of reduced precision suggests
three other ways in which sensitivity could concentrate: in certain types
of operation, such as normalisation and softmax, whose sums mixed-precision
training carries out in fp32 \citep{micikevicius2018mixed}; in a few channels
with outsized values, which dominate the quantisation error of large
language models \citep{dettmers2022llmint8}; and towards the end of the
network, if rounding errors accumulate with depth. Whether these carry
over to weather models determines whether rules from those fields can
guide reduced-precision inference here, so we test each on the same 13
dates. (i) All layers of one type -- normalisation, attention or spectral
filters, MLPs (FastNet: normalisation and linear layers) -- are kept in
fp32 with their original weights inside the fp16 network, with fp16
activations passed between them. (ii) The largest magnitude of every
output channel of every module is recorded in the fp16 forecasts. (iii)
The relative difference between the fp16 and fp32 outputs of every module
is computed during the first model step, from a fixed sample of $2^{14}$
values of each output.
FastNet runs in PyTorch's
deterministic mode in all these experiments, so that its departures
contain no run-to-run variation.

\subsection{Further diagnostics}
\label{sec:methods-further}

Latitude bands and maps (\S\ref{sec:results-where})
are computed from cos-latitude-weighted sums of the squared error, the
error and the squared departure from the fp32 forecast, accumulated on
$10^\circ\times10^\circ$ cells for every forecast, variable and lead
time; RMSE, bias and departure can then be recomputed exactly for any
union of cells.

To see whether fp16 smooths or roughens the forecast fields, we
compute, in the same runs, the zonal power spectrum of each forecast
field and of ERA5 -- the squared magnitude of its Fourier transform along
each latitude circle -- averaged with cos-latitude weights over each
$10^\circ$ latitude band. We compare the spectra of each model up to the
smallest scale it resolves, set by its patch size (Aurora and Pangu-Weather),
the wavenumbers its spectral filters act on (FourCastNet), its spherical
harmonic truncation (SFNO) or its grid spacing (FastNet).

Aurora's ensemble checkpoint is run on every fourth of the 52
dates (13) with 8 members, whose noise is set by seeds 0--7 in both fp32
and fp16, so that member $m$ of the fp16 ensemble is the same stochastic
draw as member $m$ of the fp32 ensemble. Each ensemble is scored, for
2t, msl, z500 and t850, by the
fair CRPS \citep{ferro2014fair}, the RMSE of the ensemble mean, the spread
(square root of the area-mean ensemble variance) and the fraction of the
area where ERA5 lies outside the range of the members; the precision effect
is the RMS difference between same-seed fp16 and fp32 members.

To test whether fp16's numeric range suffices, the fp16 forecasts of
all 52 dates are repeated with forward hooks on every module of each
network (not only the innermost layers, so that sums formed inside a
block are seen at its output). For every fp16 tensor a module returns,
they record the largest finite magnitude, the number of infinite or
undefined values, and the number of nonzero values below
$6.1\times10^{-5}$, both per module and per model step. Values computed
and used within a module without being returned, such as attention
scores, are not seen directly, but an overflow there would make the
module's output infinite or undefined.

For a qualitative view of high-impact weather, we track four named
storms of 2024 whose initialisation is one of the 52 dates: the
extratropical storms Isha (initialised 20 January 00\,UTC;
\citealp{metoffice2024isha}) and Kathleen (5 April 18\,UTC;
\citealp{metoffice2024kathleen}), and the hurricanes Beryl (1 July 00\,UTC;
\citealp{beven2025beryl}) and Milton (7 October 00\,UTC;
\citealp{beven2025milton}). Each storm is followed in ERA5 and in each forecast
by the minimum of mean sea-level pressure within 700\,km (extratropical)
or 400\,km (tropical) of its previous position per 6 hours, starting from
its ERA5 position (for Isha, from 21 January 12\,UTC, when it has become a
separate low). We compare track position, central pressure and the
maximum 10\,m wind within 300\,km of the centre with ERA5, for fp32 and
fp16. ERA5 underestimates the
intensity of tropical cyclones, so intensities are compared with ERA5,
not with observed intensities.

The cost of each model and precision is measured on the GPU used
throughout, an NVIDIA GeForce RTX~5080 with 16\,GB of memory, with
nothing else running on it. Each model is stepped forward from one
initial date, restarting the 10-day forecast as often as needed, until at
least 60\,s and at least 8 steps have been timed; the first step of each
forecast, which also loads the inputs, is not timed. We report the mean
time per step, the time for a 10-day forecast, the peak GPU memory
allocated by PyTorch, and the energy of a 10-day forecast: the mean board
power of the GPU, sampled every 50\,ms within the timed steps, times
their duration.

To test whether 8-bit formats are a further option, we simulate
8-bit quantisation inside the fp16 network, with the state in fp32, on
the 13 sensitivity dates over 10 days. In the weight-only configurations,
every weight of a matrix product -- the linear and convolution layers and
FourCastNet's and SFNO's spectral filters (for SFNO, the real and
imaginary parts) -- is rounded to int8 (symmetric, 127 levels) or to fp8
(e4m3) with one scale per output channel, and the network then computes
in fp16. In the configurations with 8-bit inputs, the inputs of every
linear and convolution layer are also rounded, to int8 with one scale per
token or to fp8 with one scale per tensor. Biases, normalisation layers,
embeddings and Pangu-Weather's position-dependent attention bias, which
holds 91\% of its parameters, stay in fp16, as is usual. This is the
simplest form of 8-bit quantisation, round-to-nearest without
calibration or the treatment of outliers developed for language models
\citep{dettmers2022llmint8, xiao2022smoothquant, lin2023awq}; it
reproduces the accuracy of 8-bit inference but not its speed or memory,
which we do not measure.
